\section{Review of experimental settings in PEFT studies}\label{app:peft-review}

\paragraph{Papers and comparisons.}
We review method papers associated with Hugging Face PEFT v0.21.0 \citep{mangrulkar2022peft}, using paper versions available by 26 September 2026.
Multiple supported entries from the same paper are counted once, and background citations are excluded.
The inventory contains 71 entries and 70 papers, with one entry lacking an identified method paper.
Sixty-four papers have eligible main comparisons.
The other six concern training from scratch, contributions without adaptation training or results without a trained principal comparator outside auxiliary analyses.

We read the papers, appendices, available supplements and relevant official experimental descriptions.
A comparison pairs the proposed method with its principal trained baseline for a stated pretrained model and training task.
Vanilla LoRA is preferred when present.
Several evaluations of the same adapted model remain one comparison, while separately trained tasks are distinguished.
Additional ablations do not establish controls for the main comparisons.
Capacity strata are retained within each comparison.

\paragraph{Controls and evidence.}
\Cref{tab:peft-review-criteria} defines the four controls.
Both methods must have the required evidence in the same comparison.
Comparable conditions are assessed separately for checkpoint, data, exposure, adapted modules and trainable budget.
Shared stopping rules can establish comparable exposure.
Parameter counts within 10\% of the larger count, an explicit matched design or a controlled comparison across budgets establish budget comparability.
Equal rank alone is insufficient.
Intrinsic module differences and full fine-tuning budgets are treated as inapplicable where direct matching is inappropriate.
A different budget can still support a useful efficiency comparison.

Training runs may vary initialization, data order or sampled training examples.
Their uncertainty is distinguished from variation across generated samples, subjects, evaluation subsets or evaluation resamples.
Undefined error bars and dispersion pooled across training seeds and other units do not establish uncertainty attributable specifically to training.
A numerical search space does not establish equal tuning effort or separate validation selection.
An explicitly exhaustive design with one candidate can document disclosure.
Selection using the reported validation split is recorded separately from selection using an independent split.

Two AI-assisted source reviews were compared, and disagreements were revisited against the primary sources.
Differences in eligibility, reviewed versions, baseline attribution and experimental scope were reconciled using the definitions above.
Ambiguous or contradictory evidence remains unclear.
Relevant official descriptions provide supporting evidence, while generic defaults do not establish which procedures produced reported results.
AI-assisted review can retain correlated errors.

\paragraph{Aggregation and coverage.}
A paper documents a control if at least one eligible main comparison satisfies it.
An explicit departure requires evidence that every eligible comparison fails the definition.
Otherwise, the classification is unclear.
Inapplicable cases are excluded from the relevant denominator.
Missing information is never interpreted as evidence that a procedure was not performed.
For incomplete comparison coverage, known positives are retained, while universal claims are withheld unless a statement explicitly covers all main results.

\Cref{tab:peft-review-coverage,tab:peft-review-domains} report coverage and domain summaries.
Sixteen papers have unresolved comparison coverage.
No reviewed comparison establishes all four controls together under these definitions.
Requiring three training runs retains the same nine replication positives, while requiring explicitly stated training randomness retains seven.
Excluding learning-rate selection on the reported or an unnamed split leaves four selection positives.

\paragraph{Scope of interpretation.}
Library membership defines a practical sample spanning several domains and research objectives.
It is not a representative sample of all PEFT research.
Domains overlap at paper level, and controls documented in different comparisons cannot be combined into a single qualifying experiment.
The cohort and validation-selection definition differ from \citet{lee2026learningrate}, so the percentages cannot be compared directly.
The review characterizes available evidence and does not estimate how undocumented procedures would change method rankings.

\begin{table}[!htbp]
\captionsetup{labelsep=period}
\caption{Definitions used in the review. Evidence must cover both methods and the stated main comparison. These are operational criteria for this review.}
\label{tab:peft-review-criteria}
\centering
\begin{tabular}{@{}p{.25\linewidth}p{.71\linewidth}@{}}
\toprule
Control & Required evidence \\
\midrule
Training runs and uncertainty & At least two independent training runs per method with a defined SD, SE or CI attributable to training variation. \\
LR selection & At least three learning rates evaluated per method, with separate selection using a stated validation criterion. \\
LR candidates & Numerical candidate values or a search range stated for both methods. \\
Comparable conditions & Shared checkpoint and data, comparable exposure, and adapted modules and trainable budgets matched or addressed through an appropriate controlled comparison. \\
\bottomrule
\end{tabular}
\end{table}

\begin{table}[!htbp]
\captionsetup{labelsep=period}
\caption{Paper-level coverage of experimental controls. The first three columns partition the 64 eligible papers using the at least one comparison rule. The final column requires established coverage of every applicable main comparison. Unresolved coverage prevents a universal positive.}
\label{tab:peft-review-coverage}
\centering
\begin{tabular*}{\linewidth}{@{\extracolsep{\fill}}lrrrr@{}}
\toprule
Control & Documented & Explicit departure & Unclear & \shortstack{All main\\comparisons} \\
\midrule
Training runs and uncertainty & 9 & 5 & 50 & 2 \\
LR selection & 7 & 4 & 53 & 1 \\
LR candidates & 13 & 0 & 51 & 4 \\
Comparable conditions & 25 & 10 & 29 & 1 \\
\bottomrule
\end{tabular*}
\end{table}

\begin{table}[!htbp]
\captionsetup{labelsep=period}
\caption{Documented controls within each domain. Entries give qualifying papers over applicable papers. A paper can occur in several domains, so rows cannot be summed. The decoder subset restricts comparisons to standard decoder adaptation against vanilla LoRA, with the exclusions described by Lee et al. Its definitions and sampling frame still differ from theirs.}
\label{tab:peft-review-domains}
\centering
\begin{tabular*}{\linewidth}{@{\extracolsep{\fill}}lrrrr@{}}
\toprule
Domain & \shortstack{Training\\uncertainty} & LR selection & LR candidates & Conditions \\
\midrule
Decoder language & 4/51 & 4/51 & 7/51 & 18/51 \\
Other language & 5/40 & 3/40 & 6/40 & 7/40 \\
Vision and multimodal & 3/20 & 1/20 & 3/20 & 2/20 \\
Diffusion & 0/11 & 0/11 & 1/11 & 3/11 \\
Other & 0/1 & 0/1 & 0/1 & 0/1 \\
\midrule
Decoder subset & 3/40 & 2/40 & 4/40 & 12/40 \\
\bottomrule
\end{tabular*}
\end{table}

